\Needspace{6\baselineskip}\Article{331}
In Abschnitt IV haben wir ausführlich gezeigt, wie man die Primzahlen, deren Rest eine gegebene Zahl \(r\) ist (die wir als durch kein Quadrat teilbar voraussetzen dürfen), von denen unterscheidet, deren Nichtrest sie ist, oder die Teiler des Ausdrucks \(xx-r\) von den Nichtteilern: Alle ersteren sind unter gewissen Formeln der Art \(rz+a,\allowbreak \ rz+b\) usw. oder der Art \(4rz+a,\allowbreak \ 4rz+b\) usw. enthalten, die letzteren unter anderen ähnlichen Formeln. Wenn \(r\) eine sehr kleine Zahl ist, lassen sich die Ausschlüsse mithilfe dieser Formeln
\SourceLeaf{409}
sehr bequem durchführen; zum Beispiel sind alle Zahlen der Form \(4z+3\) auszuschließen, wenn \(r=-1\) gilt; alle Zahlen der Formen \(8z+3\) und \(8z+5\), wenn \(z=2\) gilt, usw. Doch können wir nicht immer solche Reste der vorgelegten Zahl \(M\) finden, und für einen großen Wert von \(r\) ist die Anwendung der Formeln nicht bequem genug. Es ist daher ein gewaltiger Vorteil und erleichtert die Ausschlussarbeit außerordentlich, wenn für eine hinreichend große Menge positiver und negativer, durch kein Quadrat teilbarer Zahlen \((r)\) bereits eine \textit{Tafel} vorliegt, in der die Primzahlen, deren Reste die einzelnen \((r)\) sind, von denen unterschieden werden, deren Nichtreste sie sind. Eine solche Tafel lässt sich ebenso einrichten wie die diesem Werk am Ende beigefügte und oben schon beschriebene Probe; damit sie aber für den gegenwärtigen Zweck hinreichend großen Nutzen bietet, müssen die am Rand stehenden Primzahlen (die Moduln) viel weiter, etwa mindestens bis \(1000\) oder \(10000\), fortgesetzt werden. Außerdem wird die Anwendung erheblich bequemer, wenn unter den Überschriften auch zusammengesetzte und negative Zahlen aufgenommen werden, obwohl dies, wie Abschnitt IV zeigt, nicht unbedingt nötig ist. Am bequemsten wird der Gebrauch einer solchen Tafel, wenn man ihre einzelnen senkrechten Spalten ausschneidet und auf Streifen oder Stäbe (ähnlich den Neperschen) klebt, sodass die jeweils benötigten, nämlich diejenigen zu den Zahlen \(r,\allowbreak r',\allowbreak r''\) usw., den Resten der vorgelegten, in Faktoren zu zerlegenden Zahl, einzeln untersucht werden können. Legt man sie \textit{richtig} neben die erste Spalte der Tafel (welche die Moduln enthält), das heißt so, dass die Stellen der einzelnen Stäbe, die derselben Primzahl der ersten Spalte entsprechen, mit ihr in gerader Richtung oder auf derselben waagerechten Linie liegen, so erkennt man offenbar auf einen Blick die Primzahlen, die nach den Ausschlüssen mit den Resten \(r,\allowbreak r',\allowbreak r''\) in \(\Omega\) übrig bleiben: Es sind gerade diejenigen Zahlen der ersten Spalte, neben denen sich auf \textit{allen} benachbarten Stäben ein Strich befindet; alle Zahlen, neben denen auf \textit{irgendeinem} Stab eine leere Stelle liegt, sind zu verwerfen. Ein Beispiel wird dies hinreichend verdeutlichen. Weiß man auf irgendeine Weise, dass die Zahlen \(-6,\allowbreak +13,\allowbreak -14,\allowbreak +17,\allowbreak +37,\allowbreak -53\) Reste von \(997331\) sind, so sind die erste Spalte (die in diesem Fall bis \(997\), der Primzahl unmittelbar unter \(\surd{997331}\), reichen muss) und die Streifen mit den Überschriften \(-6,\allowbreak +13\) usw. zusammenzulegen. Hier ist ein Teil des so entstehenden Schemas:
\SourceLeaf{410}
\[\renewcommand{\arraystretch}{1.05}\begin{array}{|r|c|c|c|c|c|c|}&-6&+13&-14&+17&+37&-53\\\hline 3&-&-&-&&-&-\\5&-&&-&&&\\7&-&&-&&-&\\11&-&&&&-&\\13&&-&-&-&&-\\17&&-&&-&&-\\19&&&-&-&&-\\23&&-&-&&&-\\&&&\text{usw.}&&&\\113&&-&-&&&-\\127&-&-&-&-&-&-\\131&-&-&-&&&\\&&&\text{usw.}&&&\end{array}\]

Wie man hier auf einen Blick erkennt, dass von den Primzahlen, \textit{die in diesem Teil des Schemas enthalten sind}, allein \(127\) nach den Ausschlüssen mit den Resten \(-6,\allowbreak 13\) usw. in \(\Omega\) übrig bleibt, so zeigt das vollständige, bis \(997\) fortgesetzte Schema, dass \textit{überhaupt} keine andere in \(\Omega\) übrig bleibt; versucht man die Division, so findet man, dass \(997331\) tatsächlich durch \(127\) teilbar ist. Auf diese Weise ist jene Zahl also in die Primfaktoren \(127\times7853\) zerlegt\HistoricalNote{*)}{Der Autor würde den recht umfangreichen Apparat der hier beschriebenen Tafel, den er für seinen eigenen Gebrauch anfertigen ließ, gern veröffentlichen, wenn die geringe Zahl derer, denen er nützlich sein kann, zur Deckung der Kosten eines solchen Unternehmens ausreichte. Möchte sich inzwischen ein Freund der Arithmetik, nachdem er die Grundsätze gründlich verstanden hat, aus eigener Kraft eine solche Tafel anlegen, so wird es dem Autor große Freude machen, ihm alle Hilfsmittel und Kunstgriffe brieflich mitzuteilen.}.

Aus dieser Darstellung geht außerdem hinreichend hervor, dass Reste, die nicht zu groß sind oder sich wenigstens in nicht zu große Primfaktoren zerlegen lassen, besonders nützlich sind: Der unmittelbare Gebrauch der Hilfstafel reicht nicht über die Zahlen in ihren Überschriften hinaus, und der mittelbare umfasst nur solche, die sich in Faktoren zerlegen lassen, welche in der Tafel enthalten sind.

\Needspace{6\baselineskip}\Article{332}
Wir werden drei verschiedene Verfahren angeben, um Reste einer gegebenen Zahl \(M\) zu finden. Ihrer Darstellung schicken wir zwei Beobachtungen voraus, mit deren Hilfe man aus weniger geeigneten Resten einfachere ableiten kann. \textit{Erstens}: Ist eine Zahl \(akk\), die durch das Quadrat \(kk\) teilbar ist (das als zu \(M\) teilerfremd vorausgesetzt wird), ein Rest von \(M\), so ist auch \(a\) ein Rest; deshalb sind Reste, die durch große
\SourceLeaf{411}
Quadrate teilbar sind, ebenso nützlich wie kleine, und wir werden voraussetzen, dass aus allen Resten, welche die folgenden Verfahren liefern, die quadratischen Faktoren sofort entfernt worden sind. \textit{Zweitens}: Sind zwei oder mehr Zahlen Reste, so ist auch ihr Produkt ein Rest. Verbindet man diese Beobachtung mit der vorigen, so lässt sich sehr häufig aus mehreren Resten, die nicht alle einfach genug sind, ein anderer sehr einfacher ableiten, sofern jene viele gemeinsame Faktoren enthalten. Aus diesem Grund sind auch solche Reste sehr geeignet, die aus vielen nicht zu großen Faktoren bestehen, und es wird zweckmäßig sein, sie alle sofort in ihre Faktoren zu zerlegen. Die Bedeutung dieser Beobachtungen wird durch Beispiele und häufigen Gebrauch besser verständlich als durch Regeln.

I. Das einfachste und für diejenigen, die durch häufige Übung schon einige Gewandtheit erworben haben, bequemste Verfahren besteht darin, \(M\) oder allgemeiner irgendein Vielfaches von \(M\) auf beliebige Weise in zwei Teile \(kM=a+b\) zu zerlegen (seien beide positiv oder einer positiv und der andere negativ). Ihr Produkt mit umgekehrtem Vorzeichen ist dann ein Rest von \(M\); denn es ist \(-ab\equiv aa\equiv bb\ (\mathrm{mod.}\ M)\) und daher \(-ab\,R\,M\). Die Zahlen \(a,\allowbreak b\) sind so zu wählen, dass ihr Produkt durch ein großes Quadrat teilbar und der Quotient entweder klein oder wenigstens in nicht zu große Faktoren zerlegbar ist; dies lässt sich immer ohne Schwierigkeit erreichen. Besonders empfiehlt es sich, für \(a\) ein Quadrat oder das Doppelte oder Dreifache eines Quadrats usw. zu nehmen, das von der Zahl \(M\) um eine entweder kleine oder in bequeme Faktoren zerlegbare Zahl abweicht. So findet man zum Beispiel \[\begin{aligned}
997331&=999^2-2.5.67\\
&=994^2+5.11.13^2=2.706^2+3.17.3^2\\
&=3.575^2+11.31.4^2=3.577^2-7.13.4^2\\
&=3.578^2-7.19.37=11.299^2+2.3.5.29.4^2\\
&=11.301^2+5.12^2
\end{aligned}\] usw. Daraus erhält man die folgenden Reste: \(2.5.67,\allowbreak \ -5.11,\allowbreak \ -2.3.17,\allowbreak \ -3.11.31,\allowbreak \ 3.7.13,\allowbreak \ 3.7.19.37,\allowbreak -2.3.5.11.29\); die letzte Zerlegung liefert den Rest \(-5.11\), den wir schon haben. Anstelle der Reste \(-3.11.31,\allowbreak \ -2.3.5.11.29\) können wir \(3.5.31,\allowbreak \ 2.3.29\) nehmen, die aus ihrer Verbindung mit \(-5.11\) entstehen.

\begin{samepage}
II. Das zweite und dritte Verfahren beruhen darauf, dass, wenn zwei binäre Formen \((A,\allowbreak B,\allowbreak C),\allowbreak \ (A',\allowbreak B',\allowbreak C')\) derselben Determinante \(M\) oder \(-M\) oder allgemeiner \(\pm kM\) zum selben Geschlecht gehören, die Zahlen \(AA',\allowbreak \ AC',\allowbreak \ A'C\) Reste von \(kM\) sind. Dies erkennt man ohne Mühe daran, dass jede charakteristische Zahl der einen Form, etwa \(m\), auch eine charakteristische Zahl der anderen ist und somit \(mA,\allowbreak \ mC,\allowbreak \)
\SourceLeaf{412}
\(mA',\allowbreak \ mC'\) sämtlich Reste von \(kM\) sind. Ist also \((a,\allowbreak b,\allowbreak a')\) eine reduzierte Form mit positiver Determinante \(M\) oder allgemeiner \(kM\), und sind \((a',\allowbreak b',\allowbreak a''),\allowbreak \ (a'',\allowbreak b'',\allowbreak a''')\) usw. Formen aus ihrer Periode, also ihr äquivalent und erst recht im selben Geschlecht enthalten, so sind die Zahlen \(aa',\allowbreak \ aa'',\allowbreak \ aa'''\) usw. sämtlich Reste von \(M\). Eine große Anzahl von Formen einer solchen Periode lässt sich mithilfe des Algorithmus aus Artikel 187 sehr leicht berechnen; die einfachsten Reste ergeben sich meistens, wenn man \(a=1\) setzt; diejenigen, die zu große Faktoren enthalten, sind zu verwerfen. Hier sind die Anfänge der Perioden der Formen \((1,\allowbreak 998,\allowbreak -1327)\) und \((1,\allowbreak 1412,\allowbreak -918)\), deren Determinanten \(997331,\allowbreak \ 1994662\) sind: \[\begin{array}{rrr|rrr}\bigl(1,&998,&-1327\bigr)&\bigl(1,&1412,&-918\bigr)\\ \bigl(-1327,&329,&670\bigr)&\bigl(-918,&1342,&211\bigr)\\ \bigl(670,&341,&-1315\bigr)&\bigl(211,&1401,&-151\bigr)\\ \bigl(-1315,&974,&37\bigr)&\bigl(-151,&1317,&1723\bigr)\\ \bigl(37,&987,&-626\bigr)&\bigl(1723,&406,&-1062\bigr)\\ \bigl(-626,&891,&325\bigr)&\bigl(-1062,&656,&1473\bigr)\\ \bigl(325,&734,&-1411\bigr)&\bigl(1473,&817,&-901\bigr)\\ \bigl(-1411,&677,&382\bigr)&\bigl(-901,&985,&1137\bigr)\\ \bigl(382,&851,&-715\bigr)&\multicolumn{3}{c}{\mathrm{usw.}}\end{array}\] Es sind also Reste der Zahl \(997331\) sämtliche Zahlen \(-1327,\allowbreak 670\) usw.; lässt man diejenigen mit zu großen Faktoren beiseite, so erhält man die folgenden: \(2.5.67,\allowbreak \ 37,\allowbreak \ 13,\allowbreak -17.83,\allowbreak \ -5.11.13,\allowbreak \ -2.3.17,\allowbreak \ -2.59,\allowbreak \ -17.53\). Den Rest \(2.5.67\) sowie \(-5.11\), der aus der Verbindung des dritten mit dem fünften entsteht, hatten wir schon oben gefunden.
\par\end{samepage}

III. Ist \(C\) irgendeine von der Hauptklasse verschiedene Klasse von Formen mit negativer Determinante \(-M\) oder allgemeiner \(-kM\) und ihre Periode \(2C,\allowbreak \ 3C\) usw. (Artikel 307), so gehören die Klassen \(2C,\allowbreak \ 4C\) usw. zum Hauptgeschlecht, dagegen \(3C,\allowbreak \ 5C\) usw. zum selben Geschlecht wie \(C\). Ist also \((a,\allowbreak b,\allowbreak c)\) eine (einfachste) Form aus \(C\) und \((a',\allowbreak b',\allowbreak c')\) eine Form aus irgendeiner Klasse jener Periode, etwa aus \(nC\), so ist entweder \(a'\) oder \(aa'\) ein Rest von \(M\), je nachdem, ob \(n\) gerade oder ungerade ist (im ersten Fall offenbar auch \(c'\), im zweiten \(ac',\allowbreak \ ca'\) und \(cc'\)). Die Periode, das heißt die einfachsten Formen in ihren Klassen, lässt sich mit erstaunlicher Leichtigkeit entwickeln, wenn \(a\) sehr klein ist, insbesondere im Fall \(=3\), was sich immer erreichen lässt, wenn \(kM\equiv2\ (\mathrm{mod.}\ 3)\) gilt. Hier ist der Anfang der Periode der Klasse, welche die Form \((3,\allowbreak 1,\allowbreak 332444)\) enthält:
\SourceLeaf{413}
\[\begin{array}{crrr|crrr}C&\bigl(3,&1,&332444\bigr)&6C&\bigl(729.&-209,&1428\bigr)\\ 2C&\bigl(9,&-2,&110815\bigr)&7C&\bigl(476,&209,&2187\bigr)\\ 3C&\bigl(27,&7,&36940\bigr)&8C&\bigl(1027,&342,&1085\bigr)\\ 4C&\bigl(81,&34,&12327\bigr)&9C&\bigl(932,&-437,&1275\bigr)\\ 5C&\bigl(243,&34,&4109\bigr)&10C&\bigl(425,&12,&2347\bigr)\end{array}\] Daraus gehen die Reste (unter Verwerfung der unbrauchbaren) \(3.476,\allowbreak \ 1027,\allowbreak \ 1085,\allowbreak \ 425\) hervor, oder (nach Entfernung der quadratischen Faktoren) \(3.7.17,\allowbreak \ 13.79,\allowbreak \ 5.7.31,\allowbreak \ 17\). Verbindet man diese passend mit den acht in II gefundenen Resten, so erhält man leicht die folgenden zwölf: \(-2.3,\allowbreak \ 13,\allowbreak \ -2.7,\allowbreak 17,\allowbreak \ 37,\allowbreak \ -53,\allowbreak \ -5.11,\allowbreak \ 79,\allowbreak \ -83,\allowbreak \ -2.59,\allowbreak \ -2.5.31,\allowbreak \ 2.5.67\); die ersten sechs sind dieselben, die wir in Artikel 331 verwendet haben. Die Reste \(19\) und \(-29\) hätten hinzugefügt werden können, wenn wir auch die in I gefundenen hätten verwenden wollen; die übrigen dort erhaltenen hängen bereits von denen ab, die wir hier entwickelt haben.

\Needspace{6\baselineskip}\Article{333}
Das \textsc{zweite Verfahren}, eine gegebene Zahl \(M\) in Faktoren zu zerlegen, geht von der Betrachtung der Werte eines Ausdrucks wie \(\surd{-D}\ (\mathrm{mod.}\ M)\) aus und stützt sich auf die folgenden Beobachtungen.

I. Ist \(M\) eine Primzahl oder eine Potenz einer Primzahl (wobei die Primzahl ungerade ist und \(D\) nicht teilt), so ist \(-D\) ein Rest oder Nichtrest von \(M\), je nachdem, ob \(M\) unter einer Form der Teiler oder einer Form der Nichtteiler von \(xx+D\) enthalten ist; im ersten Fall hat der Ausdruck \(\surd{-D}\ (\mathrm{mod.}\ M)\) nur zwei verschiedene Werte, die entgegengesetzt sind.

II. Ist dagegen \(M\) zusammengesetzt, etwa \(=pp'p''\) usw., wobei \(p,\allowbreak p',\allowbreak p''\) usw. Primzahlen (verschiedene, ungerade, die \(D\) nicht teilen) oder Potenzen solcher Primzahlen bezeichnen, so ist \(-D\) nur dann ein Rest von \(M\), wenn es ein Rest der einzelnen \(p,\allowbreak p',\allowbreak p''\) usw. ist, das heißt, wenn diese Zahlen alle unter Formen der Teiler von \(xx+D\) enthalten sind. Bezeichnet man die Werte des Ausdrucks \(\surd{-D}\) nach den Moduln \(p,\allowbreak p',\allowbreak p''\) usw. jeweils mit \(\pm r,\allowbreak \ \pm r',\allowbreak \ \pm r''\) usw., so entstehen sämtliche Werte desselben Ausdrucks nach dem Modul \(M\), indem man die Zahlen ermittelt, die nach \(p\) entweder \(\equiv r\) oder \(\equiv-r\), nach \(p'\) entweder \(\equiv r'\) oder \(\equiv-r'\) usw. sind. Ihre Anzahl ist daher \(=2^\mu\), wobei \(\mu\) die Anzahl der Zahlen \(p,\allowbreak p',\allowbreak p''\) usw. bezeichnet. Sind diese Werte also \(R,\allowbreak -R,\allowbreak R',\allowbreak -R',\allowbreak R''\) usw., so gilt von selbst \(R\equiv R\) nach allen \(p,\allowbreak p',\allowbreak p''\) usw., aber nach keiner von ihnen \(R\equiv-R\). Daher ist der größte gemeinsame
\SourceLeaf{414}
Teiler von \(M\) und \(R-R\) gleich \(M\), und \(1\) ist der größte gemeinsame Teiler von \(M\) und \(R+R\). Zwei Werte aber, die weder identisch noch entgegengesetzt sind, wie \(R\) und \(R'\), müssen nach einer oder mehreren der Zahlen \(p,\allowbreak p',\allowbreak p''\) usw., jedoch nicht nach allen, kongruent sein, und nach den übrigen gilt \(R\equiv-R'\); folglich ist das Produkt der ersteren der größte gemeinsame Teiler von \(M\) und \(R-R'\) und das Produkt der letzteren der größte gemeinsame Teiler von \(M\) und \(R+R'\). Daraus folgt leicht: Berechnet man alle größten gemeinsamen Teiler von \(M\) und den Differenzen zwischen den einzelnen Werten des Ausdrucks \(\surd{-D}\ (\mathrm{mod.}\ M)\) und einem beliebig gegebenen Wert, so enthält ihre Menge die Zahlen \(1,\allowbreak p,\allowbreak p',\allowbreak p''\) usw. und alle Produkte aus je zweien, dreien usw. dieser Zahlen. \textit{Auf diese Weise lassen sich also aus} \textit{den Werten jenes Ausdrucks die Zahlen \(p,\allowbreak p',\allowbreak p''\) usw. ermitteln.}

Da ferner das Verfahren aus Artikel 327 die einzelnen Werte auf die Werte von Ausdrücken der Form \(\frac mn\ (\mathrm{mod.}\ M)\) zurückführt, wobei der Nenner \(n\) zu \(M\) teilerfremd ist, brauchen diese Ausdrücke selbst für den gegenwärtigen Zweck nicht einmal berechnet zu werden. Denn der größte gemeinsame Teiler von \(M\) und der Differenz zwischen \(R\) und \(R'\), die mit \(\frac mn,\allowbreak \frac{m'}{n'}\) übereinstimmen, ist offenbar auch der größte gemeinsame Teiler von \(M\) und \(nn'(R-R')\). Oder, was dasselbe ist, von \(M\) und \(mn'-m'n\), zu dem \(nn'(R-R')\) nach dem Modul \(M\) kongruent ist.

\Needspace{12\baselineskip}\Article{334}
Die vorstehenden Beobachtungen lassen sich auf unser Problem auf zwei Arten anwenden. Die erste entscheidet nicht nur, ob die vorgelegte Zahl \(M\) eine Primzahl oder zusammengesetzt ist, sondern liefert im letzteren Fall auch die Faktoren selbst; die zweite hat den Vorteil, dass sie meist eine raschere Rechnung ermöglicht, liefert aber manchmal die Faktoren der zusammengesetzten Zahlen, die auch sie vollständig von den Primzahlen unterscheidet, erst nach mehrmaliger Wiederholung.

I. Man suche eine negative Zahl \(-D\), die ein quadratischer Rest von \(M\) ist; dazu können die in Artikel 332 unter I und II angegebenen Verfahren dienen. An sich ist die Wahl des Restes beliebig, und anders als beim vorigen Verfahren braucht \(D\) hier keine kleine Zahl zu sein; doch wird die Rechnung umso kürzer, je kleiner die Anzahl der Klassen binärer Formen in den einzelnen eigentlich primitiven Geschlechtern mit der Determinante \(-D\) ist. Deshalb sind besonders solche Reste geeignet, die unter den \(65\) Zahlen aus Artikel 303 enthalten sind, falls sich solche anbieten. So wäre für \(M=997331\) von allen oben gefundenen negativen Resten \(-102\) am besten
\SourceLeaf{415}
geeignet. Man ermittle alle verschiedenen Werte des Ausdrucks \(\surd{-D}\ (\mathrm{mod.}\ M)\); ergeben sich nur zwei (entgegengesetzte), so ist \(M\) sicher eine Primzahl oder eine Potenz einer Primzahl; sind es mehr, etwa \(2^\mu\), so ist \(M\) aus \(\mu\) verschiedenen Primzahlen oder Potenzen von Primzahlen zusammengesetzt, die man als Faktoren mit dem Verfahren des vorigen Artikels ermitteln kann. Ob diese Faktoren Primzahlen oder Potenzen von Primzahlen sind, lässt sich an sich sehr leicht erkennen; außerdem zeigt schon das Verfahren, durch das man die Werte des Ausdrucks \(\surd{-D}\) findet, von selbst alle Primzahlen an, von denen irgendeine Potenz \(M\) teilt. Ist nämlich \(M\) durch das Quadrat einer Primzahl \(\pi\) teilbar, so hat jene Rechnung sicher auch eine oder mehrere Darstellungen der Zahl \(M\), \(M=amm+2bmn+cnn\), geliefert, in denen der größte gemeinsame Teiler der Zahlen \(m,\allowbreak n\) gleich \(\pi\) ist (und zwar deshalb, weil in diesem Fall \(-D\) auch ein Rest von \(\frac{M}{\pi\pi}\) ist). Ist dagegen keine Darstellung entstanden, in der \(m\) und \(n\) einen gemeinsamen Teiler haben, so ist dies ein sicheres Zeichen dafür, dass \(M\) durch kein Quadrat teilbar ist und daher alle \(p,\allowbreak p',\allowbreak p''\) usw. Primzahlen sind.

Beispiel. Mit dem oben angegebenen Verfahren findet man vier Werte des Ausdrucks \(\surd{-408}\ (\mathrm{mod.}\ 997331)\), die mit den Werten von \(\pm\frac{1664}{113},\allowbreak \ \pm\frac{2824}{3}\) übereinstimmen; die größten gemeinsamen Teiler von \(997331\) mit \(3.1664-113.2824\) und \(3.1664+113.2824\), also mit \(314120\) und \(324104\), sind \(7853\) und \(127\), woraus wie oben \(997331=127.7853\) folgt.

II. Man nehme irgendeine negative Zahl \(-D\), sodass \(M\) unter einer Form der Teiler von \(xx+D\) enthalten ist. An sich ist die Wahl einer solchen Zahl beliebig; der Bequemlichkeit halber ist aber vor allem darauf zu achten, dass die Anzahl der Klassen in den Geschlechtern mit der Determinante \(-D\) möglichst klein ist. Eine solche Zahl zu finden, bereitet übrigens keine Schwierigkeit, wenn man durch Probieren vorgeht; denn gewöhnlich ist unter einer beträchtlichen Menge versuchter Zahlen \(M\) für ungefähr ebenso viele unter einer Form der Teiler wie unter einer Form der Nichtteiler enthalten. Es wird daher am zweckmäßigsten sein, den Versuch mit den \(65\) Zahlen aus Artikel 303 (und zwar mit den größten) zu beginnen und, falls sich keine als geeignet erweisen sollte (was jedoch im Allgemeinen nur einmal unter \(16384\) Fällen eintritt), zu anderen überzugehen, bei denen die einzelnen Geschlechter je zwei Klassen enthalten. — Dann ermittle man die Werte des Ausdrucks \(\surd{-D}\ (\mathrm{mod.}\ M)\) und leite, wenn sich welche finden, daraus die Faktoren von \(M\) auf genau dieselbe Weise wie oben ab; ergeben sich dagegen keine Werte, sodass \(-D\) ein Nichtrest von \(M\) ist, so kann \(M\) sicher weder eine Primzahl noch eine Potenz einer Primzahl sein. Wünscht man in diesem
\SourceLeaf{416}
Fall die Faktoren selbst, so muss man entweder dieselbe Operation mit anderen Werten für \(D\) wiederholen oder zu einem anderen Verfahren greifen.

So findet man zum Beispiel durch Probieren, dass \(997331\) unter einer Form der Nichtteiler von \(xx+1848,\allowbreak \ xx+1365,\allowbreak \ xx+1320\), aber unter einer Form der Teiler von \(xx+840\) enthalten ist; für die Werte des Ausdrucks \(\surd{-840}\ (\mathrm{mod.}\ 997331)\) erhält man die Ausdrücke \(\pm\frac{1272}{163},\allowbreak \pm\frac{3288}{125}\), aus denen dieselben Faktoren wie zuvor folgen. — Wer weitere Beispiele wünscht, möge Artikel 328 heranziehen: Das erste zeigt, dass \(5428681=307.17683\) gilt; das zweite, dass \(4272943\) eine Primzahl ist; das dritte, dass \(997331\) sicher aus mehreren Primzahlen zusammengesetzt ist.

\[\begin{array}{ccc}*& &*\\ &*&\end{array}\]

Der Umfang dieses Werkes hat uns hier nur erlaubt, die Hauptpunkte der beiden Verfahren zur Ermittlung der Faktoren auszuführen; eine eingehendere Untersuchung zusammen mit weiteren Hilfstafeln und anderen Hilfsmitteln behalten wir uns für eine andere Gelegenheit vor.
\BookRule{0.22\textwidth}
\clearpage
